\documentclass[11pt,reqno]{amsart}

\pdfinfoomitdate=1
\pdftrailerid{}
\pdfsuppressptexinfo=-1

\usepackage{microtype}
\microtypesetup{expansion=false}
\usepackage{mathtools}
\usepackage{amssymb}
\usepackage{url}
\usepackage[hidelinks]{hyperref}
\hypersetup{
  pdfauthor={},
  pdftitle={},
  pdfsubject={},
  pdfkeywords={},
  pdfcreator={},
  pdfproducer={}
}

\setlength{\emergencystretch}{1em}
\numberwithin{equation}{section}

\newtheorem{theorem}{Theorem}[section]
\newtheorem{lemma}[theorem]{Lemma}
\newtheorem{proposition}[theorem]{Proposition}

\newcommand{\doi}[1]{\href{https://doi.org/#1}{\nolinkurl{#1}}}

\date{July 26, 2026}

\begin{document}

\section{Proof}

\begin{theorem}[Imbalance conjecture]
Let $G$ be a finite simple graph such that $d_G(u)\ne d_G(v)$ for every
edge $uv$.  Then the multiset
\[
  \bigl(\lvert d_G(u)-d_G(v)\rvert\bigr)_{uv\in E(G)}
\]
is the degree multiset of a finite simple graph.
\end{theorem}

\begin{proof}
Write $d(v)=d_G(v)$, $D=\max_{v\in V(G)}d(v)$, and
$\iota(uv)=\lvert d(u)-d(v)\rvert$.  For $A\subseteq E(G)$ with
$\lvert A\rvert=k$, put
\[
  T_A=\sum_{e\in E(G)\setminus A}\min\{k,\iota(e)\}.
\]
We first prove
\begin{equation}\label{eq:capacity}
  T_A\ge k(D-k)^+,
\end{equation}
where $q^+=\max\{q,0\}$.

The claim is immediate when $k=0$ or $k\ge D$.  Suppose that
$0<k<D$, set $r=D-k$, and choose a vertex $x$ of degree $D$.  Let $a$
be the number of edges of $A$ incident with $x$.  Of the $D-a$ edges at
$x$ outside $A$, choose $r$ and let $R$ be their opposite endpoints; let
$X$ be the opposite endpoints of the other such edges.  Thus
$\lvert X\rvert=k-a$.

For $z\in R$, maximality of $d(x)$ and the hypothesis on adjacent
degrees give $d(z)<D$ and $\iota(xz)=D-d(z)$.  Consequently,
\begin{equation}\label{eq:reservoir}
  \min\{k,\iota(xz)\}+(d(z)-r)^+=k.
\end{equation}
Indeed, if $d(z)\le r$, the two terms are $k$ and $0$; otherwise they
are $D-d(z)$ and $d(z)-r$.

Set $\rho=\sum_{z\in R}(d(z)-r)^+$.  Call an edge an escape edge if it
has one endpoint in $R$ and the other outside $R\cup\{x\}$, and let $F$
be the set of escape edges.  For each $z\in R$, at most $r$ possible
neighbors of $z$ lie in $R\cup\{x\}$.  Since every escape edge has
exactly one endpoint in $R$, summing over $R$ gives
\begin{equation}\label{eq:escape}
  \lvert F\rvert\ge\rho.
\end{equation}

By \eqref{eq:reservoir}, the $r$ chosen edges at $x$ contribute
$rk-\rho$ to $T_A$.  No edge of $F$ is incident with $x$, while exactly
$a$ edges of $A$ are incident with $x$; hence, by \eqref{eq:escape}, at
least $(\rho-(k-a))^+$ escape edges lie outside $A$.  Each contributes at
least one to $T_A$, as does each of the $k-a$ edges from $x$ to $X$.
These three groups of edges are pairwise disjoint, so
\[
  T_A\ge rk-\rho+(k-a)+(\rho-(k-a))^+\ge rk.
\]
As $r=D-k$, this proves \eqref{eq:capacity}.

Every edge has imbalance at most $D-1$: both endpoints have positive
degree at most $D$.  It follows that every $A\subseteq E(G)$ of size $k$
satisfies
\begin{equation}\label{eq:subset}
  \sum_{e\in A}\iota(e)
  \le k(k-1)+\sum_{e\in E(G)\setminus A}\min\{k,\iota(e)\}.
\end{equation}
If $D\le k$, the left side is at most $k(D-1)\le k(k-1)$.  If $k<D$,
it is at most
\[
  k(D-1)=k(k-1)+k(D-k),
\]
and \eqref{eq:capacity} supplies the last term in \eqref{eq:subset}.

The sum of all imbalances is even.  Since $\lvert p-q\rvert\equiv p+q$
modulo two, summing first over edges and then over incidences gives
\[
  \sum_{uv\in E(G)}\lvert d(u)-d(v)\rvert
  \equiv \sum_{uv\in E(G)}(d(u)+d(v))
  =\sum_{v\in V(G)}d(v)^2.
\]
But $d(v)^2\equiv d(v)$ modulo two, and
\[
  \sum_{v\in V(G)}d(v)=2\lvert E(G)\rvert,
\]
so the imbalance sum is even.

If $E(G)$ is empty, the conclusion is immediate.  Otherwise, write the
imbalances in nonincreasing order as $s_1\ge\dots\ge s_m$, where
$m=\lvert E(G)\rvert$.  For $1\le k\le m$, choose $A$ to consist of
$k$ edges carrying $s_1,\dots,s_k$.  Then \eqref{eq:subset} says
\[
  \sum_{i=1}^{k}s_i
  \le k(k-1)+\sum_{i=k+1}^{m}\min\{k,s_i\}.
\]
Together with the even sum just proved, these are precisely the
Erd\H{o}s--Gallai conditions.  Their sufficiency shows that
$(s_1,\dots,s_m)$ is graphic.
\end{proof}

\end{document}
